\section{Описание лабораторной работы}
\label{sec:intro}

\textbf{Тема:} реализация и анализ алгоритмов обучения многослойных нейронных сетей. Применение нейронных сетей для сжатия и восстановления образов.

\textbf{Цель:} ознакомиться, проанализировать и получить навыки реализации модели линейной рециркуляционной сети для задачи сжатия графической информации.

\textbf{Вариант 2:} реализовать модель линейной рециркуляционной сети с адаптивным коэффициентом обучения и ненормированными весами~\cite{method}.

В работе строится автоассоциативная сеть, которая сжимает вектор блока изображения в скрытое представление меньшей размерности, а затем восстанавливает исходный блок. Исследуется влияние масштаба адаптивного шага, коэффициента сжатия, типа изображения и порога ошибки на число эпох обучения. Отдельно анализируется перекрытие блоков и байтовый объём представления.
